\section{Extended version of Sections 2 to 7}\label{app:J}

\emph{This appendix gives Sections 2 to 7 in the longer form of our nine-page draft of 22 October 2026, which the main text condenses; section numbers are those of the main text. Tables J.1 to J.3 give in detail what the main text condenses into Table 1, Figure 1 and Table 2 (Appendix A gives every entry of Table J.1). Two changes from the nine-page draft: in Section 6 the resubmission of 28 October and our nominations for the final test replace two placeholders, and in Sections 5.1 and 5.2 the comparisons that were not pre-registered are marked post hoc.}

\subsection*{J.2\quad The task, the labels and the public data}

\paragraph{Task.}
A submission is a function that receives a model identifier, a reasoning effort and a list of problems, and returns one robustness verdict per problem; it runs offline in the organizers' container on a single GPU, with a time limit of one hour (one hour and ten minutes since 15 October, when evaluation moved to one call per problem) [4, 5, 6]. Five models are evaluated in the main phase: Qwen3.5-4B, Skywork-OR1-Math-7B, Olmo-3-7B-Think and DeepSeek-R1-0528-Qwen3-8B in both tracks, and gpt-oss-120b in the Main track only. Since 16 September the private validation tasks have 21 items (Main) and 14 (Small); the final test will use 30 to 50 new problems from unpublished AIMO sets [5, 6].

\paragraph{Labels.}
On every one of the 379 labeled rows of the earlier public datasets, the label equals the conjunction ``no perturbation lowered accuracy'' (no detrimental permutation; permuted accuracy 1.0). A majority rule fits these rows equally well, because none has permuted accuracy strictly between 0.5 and 1; the organizers' proposal defines a majority over perturbed variants [1], and the website states that the exact criterion differs from that outline and is not published [4]. The main-phase train set (16 September; 82 rows, 71 labeled in the revision we used, extended to 94 rows on 15 October) uses thresholds on the largest accuracy drop across perturbation families: robust at most 0.10, non-robust at least 0.25, unlabeled in between, also in the validation and test splits [5]. We use the conjunction for the earlier data and report the new rule where it applies (Section 5).

\paragraph{Public data.}
Six datasets aggregate to 1,196 rows over 335 (model, problem) pairs and 174 problems; 52 pairs repeat the DeepSeek measurements under a second model name, so 283 are distinct. For DeepSeek, the only evaluated model with substantial ground truth under the earlier rule, 164 pairs are labeled: 138 declared and 26 derived by us under the conjunction (a majority rule would call 21 of these 26 robust rather than 4); every analysis on these labels inherits that choice. Ground truth for the earlier rule exists otherwise only for 10 gpt-oss-120b pairs at low reasoning effort, five of them with a declared label. Two properties of the public data matter below. In the three aggregated datasets, the field \texttt{permutation\_type} is a JSON list if and only if the item is robust, because the aggregation merges all perturbation types for robust items and keeps only the type or types with the lowest perturbed accuracy otherwise: a classifier using it is near-perfect in development and has no such field at inference (the organizers dropped it from their import script on 1 September; it remains in the datasets). And the public validation sample, \texttt{val-sample}, contains no labeled pair with base accuracy below 0.05, while 20 of the 80 robust pairs among the 335 published pairs, all of them in \texttt{aimo-interp-challenge-sample-full}, have base accuracy 0 and are robust only trivially under the conjunction.

\paragraph{Floors.}
On 21 and 14 items the binomial standard error of an accuracy near one half is 10.9 and 13.4 points. On the validation tasks introduced on 16 September the robust base rate is 10/21 and 7/14 (Section 6), so a constant scores 11/21 and 7/14 at best. A difference of one item is 4.8 and 7.1 points.

\subsection*{J.3\quad The falsification campaign}

Validation that fits on some problems and tests on others is grouped by problem. Table J.1 condenses the campaign; Appendix A gives every entry with its baseline and caveats.

\begin{center}\begin{minipage}{\linewidth}
\noindent\textbf{Table J.1.} Channels tested against ground truth (DeepSeek and the public pairs unless stated). Accuracy differences are in points against the baseline named in each pre-registration. H1, H7 and T are computed on the 335 published pairs, counting the 52 duplicated DeepSeek pairs twice; counted once, H1 is $-$7.4, H7 $-$1.4 and T +7.1, and no verdict changes.\par\smallskip
{\footnotesize\setlength{\tabcolsep}{4pt}
\begin{tabular}{L{3.7cm}L{1.3cm}L{5.5cm}L{2.3cm}}
\toprule
\textbf{Channel} & \textbf{Entry} & \textbf{Result} & \textbf{Verdict} \\
\midrule
Problem text, 32 features, grouped 5-fold & H1 & 0.675 vs 0.761 ($-$8.7) & dead \\
Robustness as a property of the problem & H2 & $-$1.3 (z = 0.88) with duplicated rows counted once; +3.4 (z = 4.04) only when counted twice & dead \\
Labels, base accuracy, difficulty order across models & H3 to H5, H7 & $-$7.3, $-$11.8, $-$25.4, $-$1.2 & dead \\
Official uncertainty features, logistic regression, grouped CV & H10 & 0.7744 vs 0.756, inside one standard error & reaches the proposal's accuracy (method substituted) \\
Cached features for other models; per-target thresholds & H12, H13 & $-$16.5; +7.3 from thresholds chosen on the scored items, $-$2.8 when forced to rank & dead; ceiling withdrawn \\
Our families; translation; self-consistency; the organizers' families regenerated (pilot) & H19, H22, H33, H32 & kappa 0.02; kappa $-$0.10; AUROC 0.502; kappa $-$0.001 & dead \\
Representation shift under perturbation; length control & H30, H31 & AUROC 0.615 (planned threshold 0.655; character count 0.705); 0.555 & dead \\
Our linear probe on the residual stream, nested CV & H36 & AUROC 0.710; about +0.02 over character count & survives, largely explained \\
The organizers' published probe artifact & H35 & AUROC 0.739 (pre-registered aggregation) and 0.709 (its own), 0.62 and 0.57 residualized on log character count; shuffled-label controls 0.68 & survives, largely explained \\
Trivial-robust rule (robust iff base accuracy < 0.05), post hoc & T & +6.0 on the 335 public pairs, +0.0 on \texttt{val-sample} & trap \\
Identical bundles re-executed (warmup, 17 items) & H18 & 11, 9, 11 and 11 correct & noise confirmed \\
\bottomrule
\end{tabular}
}
\end{minipage}\end{center}

\paragraph{Three traps a conventional validation would not catch.}
Rule T gains six points on the public pairs and nothing on \texttt{val-sample}, because the trivially robust subpopulation is absent there (Section 2). The oracle H13 appeared to show cross-model signal (+7.3 points, still +5.5 when constant predictions are excluded), but each target's threshold was chosen among all cut points on the very items it was scored on; forced to rank the items and to predict the true number of positives, it gave $-$2.8. And H2, the only warmup hypothesis about a property of the problem that had survived, owed its survival to one measurement published under two model names and counted twice: counted once, its concordance is no longer significant against a null with per-model robust rates (z = 0.88), and predicting from it scores 1.3 points below the per-model majority.

\paragraph{Length.}
Problem length recurs throughout: the character count alone orders the DeepSeek labels at AUROC 0.705, above the representation-shift statistic of H30 (0.615). Our pre-registered entry V8 (Section 6) is the rule \texttt{len(problem) >= 98}, with the threshold at the lower median of the 174 public problems and the direction fixed on the public pairs (balanced accuracy 0.652 over the 335 published pairs, 0.631 over the 283 distinct ones).

\paragraph{Probes.}
Our logistic probe on last-token hidden states of DeepSeek, with layer and regularization chosen in an inner cross-validation, reaches AUROC 0.710 out of fold and survives its permutation null (p95 0.574, with 60 permutations where 500 were pre-registered; 0.570 with the full regularization grid). The organizers' artifact survives its null as well (Table J.1), although our pre-registration predicted that both probes would die. The length controls that follow, like our analyses of the artifact (its own inference aggregation, its control groups and the directions of its probes), were not pre-registered: we added them after seeing the data. The probe's score correlates with problem length at Spearman +0.55; in a second-level cross-validation over 20 splits, character count alone gives 0.696, the probe 0.697 and the two together 0.714, a gain of about two points; inside the tercile of the longest prompts the probe is at 0.537, indistinguishable from chance. What the probe reads is not established: residualized on log character count its score gives 0.62, and log character count residualized on the probe gives 0.62 as well. The organizers' published DeepSeek probe artifact orders our labels at AUROC 0.709 under its own inference aggregation, but 0.57 once residualized on log character count, and its own randomization controls, trained on shuffled labels, give 0.68 by themselves; at the layer it uses, each seed's probe points the same way whether trained on true or shuffled labels (cosine 0.93 to 0.998). Under its own inference aggregation, then, the artifact orders the labels mostly through problem length; under the aggregation we pre-registered (AUROC 0.739) its residual on log character count is 0.62, as for our probe, which is largely redundant with length as a predictor of the label.

\subsection*{J.4\quad Re-execution noise on the evaluation workers' GPU model}

On 8 September the organizers provided a virtual machine with two NVIDIA RTX PRO 6000 Blackwell Server Edition GPUs (96 GB each), the same model as the evaluation workers. All runs in this section used the official runtime image (torch 2.12.1, transformers 5.13.0) and the unmodified official extraction and regressor, on DeepSeek.

\begin{center}\begin{minipage}{\linewidth}
\noindent\textbf{Table J.2.} Forty public problems, twenty-three executions in six configurations; generated text hashed per item. Items crossing: items whose prediction changes across the runs of a configuration (crossings grow with the number of runs; per pair of runs, 15.3 under contention, 13.3 with batch size 1, 12.7 on the second card, 11.4 by default). Deterministic algorithms: \texttt{CUBLAS\_WORKSPACE\_CONFIG=:4096:8}, \texttt{torch.use\_deterministic\_algorithms(True, warn\_only=True)} and deterministic cuDNN; compilation disabled: \texttt{TORCH\_COMPILE\_DISABLE=1}.\par\smallskip
{\footnotesize\setlength{\tabcolsep}{4pt}
\begin{tabular}{L{4.6cm}L{0.9cm}L{2.2cm}L{2.0cm}L{1.6cm}}
\toprule
\textbf{Configuration} & \textbf{Runs} & \textbf{Byte-identical items} & \textbf{Items crossing} & \textbf{Seconds per item} \\
\midrule
Default (batch 8) & 5 & 0/40 & 24/40 & 26.0 \\
Batch size 1 & 3 & 1/40 & 20/40 & 72.3 \\
Two concurrent processes on one card & 6 & 0/40 & 32/40 & 54.7 \\
Second physical card & 3 & 0/40 & 19/40 & 26.0 \\
Deterministic algorithms and cuBLAS workspace setting & 3 & 40/40 & 0 & 27.4 \\
Compilation disabled & 3 & 40/40 & 0 & 27.8 \\
\bottomrule
\end{tabular}
}
\end{minipage}\end{center}

Crossings can only grow with the number of runs, so rows with different run counts are not directly comparable; averaged over pairs of runs, 15.3 items cross under contention, 13.3 with batch size 1, 12.7 on the second card and 11.4 under the default configuration.

\paragraph{Fourteen states.}
The twenty-three executions fall into fourteen distinct numerical states. Within a state all forty generations are byte-identical; different states share few: 61 of the 91 pairs of states share no generation and 28 share one or two, while one contention pair shares 9 and two batch-1 runs share 38. The state belongs to the process, not to the card: two runs launched in the same second on the same card landed in different states, while two of three runs on the second card reproduced the first card's default state exactly. We infer, without proof, that the state is fixed when the process starts, by a choice of kernels or algorithms that we did not identify. Two pre-registered predictions failed. Batch size 1 did not make generation more repeatable: its three runs gave three different states, with 1 of 40 items byte-identical across all three, while three of the five default runs were byte-identical to one another. The configuration without compilation, added during the run, was predicted to be byte-identical across its runs and to land in one of the default states; it was byte-identical, but in a state of its own. Because only 7 of the 40 problems are robust and crossings in both directions largely cancel, accuracy moved much less than predictions (from 21 to 24 correct under the default configuration on either card, and from 24 to 28 under contention), always below the 33 of the majority constant.

\paragraph{Candidate remedies.}
In three sequential runs on one card, either deterministic algorithms with a fixed cuBLAS workspace or disabling compilation made generation byte-identical at 5 to 7\% runtime cost. Neither was tested under contention or across cards, and three identical sequential runs is also what the default configuration gave before its fourth run diverged, so these are candidates, not fixes. Determinism also does not buy correctness: the two deterministic configurations settled on two different states, predicting 15 and 12 of 40 problems robust, with no generation in common and the same prediction on only 21 of 40. They also ran on different cards (deterministic algorithms on the first, compilation disabled on the second), so configuration and card are not separated, although the default state did reproduce across cards. The published number under any deterministic setting depends on which setting was chosen.

\paragraph{The official score channel.}
The official features and regressor score were extracted three times for the 174 public problems with the same code and image (this measurement was not pre-registered and is exploratory). Two runs landed in the same state; the third, flag-free like the second, matched it on 6 of 174 scores. The scores of the two flag-free runs correlate at Spearman 0.24 over the 164 labeled problems, and in the orientation of the regressor's own decision rule the extractions give AUROC 0.487 and 0.602. On 164 problems the difference is not significant (paired bootstrap 95\% interval $-$0.015 to 0.246): the channel's ordering does not reproduce between identical runs, which is not the same as a demonstrated change of signal. Calibration moves with the state too: the shipped threshold predicts 25.6\% of problems robust in one state and 33.5\% in the other, against a base rate of 24.4\%. Since 137 of the 164 problems are in the regressor's own training set, this is a statement about the stability of its inputs, not a test on new problems.

\subsection*{J.5\quad Can the label be reproduced?}

If the label could be regenerated from the rule its public labels follow, it could be learned for the three evaluated models that had no public labels before the main-phase train set. We identified the organizers' five perturbation families from the public data (rephrase, rename, domain, distract, typos), regenerated them with gpt-oss-120b at low reasoning effort prompted with the public examples, and gated each variant with gpt-oss-120b at medium effort (the text must differ, the length ratio must lie between 0.7 and 4.0 and the solver's answer must not change; the prompts are in Appendix I). Because the public datasets carry no answers, an answer counts as correct when it matches gpt-oss-120b's answer on the original problem. On the organizers' virtual machine (E13), we labeled four models with ten variants per family in two independent draws (A and B; about 40 accepted variants per problem), ten samples on the base problem and one per variant, and re-sampled draw A's variants with a new seed (A$'$). Generation, gating and labeling ran on that machine in our own Python environment with vLLM 0.28.0 (torch 2.13.0, transformers 5.16.1, as recorded on 17 September), outside the official runtime image; variants were generated at temperature 1.0 with at most 4,096 tokens and gated at temperature 1.0 with at most 8,192 tokens.

\begin{center}\begin{minipage}{\linewidth}
\noindent\textbf{Table J.3.} Agreement of regenerated labels (Cohen's kappa; in brackets, 95\% percentile bootstrap intervals over problems, 2,000 resamples). Robust rates are over all 164 problems (the 4 without accepted variants count as not robust); kappas are over the problems labeled in both series (159 for A and B, 160 for A$'$ and for the official labels); the main-phase column uses the problems labeled in both draws under the 0.10/0.25 thresholds.\par\smallskip
{\footnotesize\setlength{\tabcolsep}{4pt}
\begin{tabular}{L{1.5cm}L{1.1cm}L{2.1cm}L{2.1cm}L{2.2cm}L{3.0cm}}
\toprule
\textbf{Model} & \textbf{Robust in draw A} & \textbf{kappa(A, B), conjunction} & \textbf{kappa(A, A$'$), same variants, new seed} & \textbf{kappa(A, B), main-phase rule} & \textbf{kappa with official labels} \\
\midrule
Skywork & 72.6\% & 0.439 [0.28, 0.59] & 0.789 [0.67, 0.89] & 0.49 [0.17, 0.76] (132 problems) & no earlier labels \\
DeepSeek & 37.8\% & 0.413 [0.26, 0.55] & 0.464* [0.33, 0.59] & 0.56 [0.36, 0.73] (100) & 0.025 [$-$0.12, 0.18]; 0.02 [$-$0.09, 0.13] under the main-phase rule (133) \\
Olmo & 65.9\% & 0.378 [0.21, 0.54] & 0.653 [0.52, 0.78] & 0.41 [0.09, 0.69] (133) & no earlier labels \\
Qwen3.5-4B & 9.1\% & 0.178 [$-$0.06, 0.39] & not run & 0.34 [0.13, 0.53] (86) & no earlier labels \\
\bottomrule
\end{tabular}
\par\smallskip\noindent (*) On DeepSeek, A$'$ also received the answer-format line of Section 5.3, so this entry mixes seed and prompt.
}
\end{minipage}\end{center}

\paragraph{5.1 More variants do not help, and our noise does not explain the gap.}
Subsampling DeepSeek's variants from 2 to 40 per problem (20 draws per size), the mean kappa(A, B) stays between 0.36 and 0.47 with no upward trend, and the mean concordance with the official label between $-$0.02 and +0.03. Correcting the observed 0.025 for our process's unreliability gives 0.039; even at the upper end of its interval the corrected value is about 0.27, so our noise could hide at most a weak agreement. Under the main-phase rule (E13b, pre-registered on the same answers with no new generation), all three hypotheses died and none of our three predictions held in full: test-retest stays below 0.60 on every model, and, post hoc, the change of rule adds only 0.01 to 0.11 when both rules are read on the same problems. The gap is not only one of severity: under the main-phase rule, on the compared problems our DeepSeek labels are robust on 71 to 75\% against the organizers' 28 to 30\%, but, in comparisons made post hoc, our robust rate is about the same whether their label is robust or not (77.5\% and 74.2\% in draw A), and the largest accuracy drops do not correlate (Spearman $-$0.09 and $-$0.01), so no stricter threshold on our side would recover theirs.

\paragraph{5.2 What drives the instability.}
On Skywork, re-sampling the same variants with a new seed leaves much higher agreement than drawing new variants (0.79 against 0.44; under the main-phase rule and post hoc, 0.94 against 0.47 on the problems labeled in all three series), so the label depends mainly on which variants are drawn. On Olmo the seed alone accounts for about half of the instability under either rule, although under the main-phase rule this rests on 5 and 9 disagreements, too few to separate the two sources. A plug-in simulation of the organizers' process from their published prediction counts gives a test-retest kappa of 0.94 (median 0.93), but it resamples answers to the same variants and never lets a robust item fail; run on our own draw A, the same procedure gives 0.64 to 0.84, against a measured 0.18 to 0.44 under new variants. It overstates by about 0.4 on every model, so the reliability of the official label remains unmeasured. The files behind the official DeepSeek labels also used a different sampling regime: on 154 of the 164 problems most variants were answered ten times each and the base between 1 and 60 times, where we answered each variant once, which by itself makes a conjunction label stricter; and our decoding (temperature 0.6, top\_p 0.95, at most 16,384 tokens) differs from theirs (temperature 1.0, up to 100,000 tokens), published after our runs [4]. What the data show is narrower than a verdict on the label: variants produced by our process measure something else. Whether any process reproduces the label is open; on 22 September the organizers wrote that they have decided not to publish the code that reproduces their labeling, to test the generalization of submitted methods [6].

\paragraph{5.3 An unspecified prompt moves the label.}
Fed the bare problem text, Qwen3.5-4B reasons for about eight thousand tokens and closes its answer in the expected format in only 5 of its 1,640 base samples in draw A. With one added instruction line requesting the format, the share of problems passing the base condition of the label (at least 6 of 10 samples agreeing with the reference answer) moved from 0 of 164 to 145 of 164 (88\%), and its robust rate from 0\% to 9.1\%; DeepSeek passes the base condition on 147 of 164 problems with or without that line. A label defined by correctness is not comparable across prompts that the benchmark leaves unspecified.

\subsection*{J.6\quad Our entries, and the evaluations the organizers requested}

\paragraph{v6, the measuring entry.}
It loads each evaluated model, extracts the official features from the model's own greedy generation and applies the official regressor (with a gpt-oss threshold fitted at run time on the ten gpt-oss rows of \texttt{val-sample}, whose only robust label is on one of three rows of \texttt{huikang-gpt-oss-120b-aimo3:low}, by its name a fine-tune of gpt-oss-120b). On the original main-phase tasks it scored 11/21 and 10/19, with runtimes of 838 s and 509 s.

\paragraph{V8, the entry that does nothing.}
Pre-registered on 13 September with conditional predictions, its entire logic is \texttt{return [len(p) >= 98 for p in problems]}: no model is loaded and each prediction depends only on its own item. On the original tasks it scored 11/21 and 10/19, identical to v6 and to the organizers' own entries on the leaderboard (0.52 and 0.53); on 21 and 19 items these ties are compatible with length carrying little signal (95\% Wilson intervals 0.32 to 0.72 and 0.32 to 0.73) without showing that it carries none, and none of these entries scored above the majority constant (at least 11/21 and 10/19). A one-line rule that ties a pipeline loading up to five models for 8.5 to 14 minutes is a measurement of the pipeline. On the replacement tasks of 16 September it predicted robust for every item and scored 10/21 and 7/14, the robust base rates.

\paragraph{The evaluations requested for reports [5].}
On the public main-phase train set (the revision of 17 September), V8 predicts robust for 65 of 71 labeled rows and scores 46/71 (0.648), below the 52/71 (0.732) of the all-robust constant (balanced accuracy 0.442). On the private validation set it scored 10/21 and 7/14 under the original batch protocol (21 September). On 15 October the organizers moved validation to one call per problem, erased the leaderboard and re-evaluated a portion of the top submissions [6]; on 28 October we resubmitted V8 unchanged (Codabench submissions 1016713, Main, and 1016730, Small), with an all-robust constant as control (1016732 and 1016735). Called once per problem through its legacy signature \texttt{are\_robust(model\_id, problems)}, V8 again predicted robust for every item, as its logged counts show, and scored 10/21 and 7/14, the robust base rates; the constant scored the same. For the final test we nominated V8 (1016713) and the all-robust constant (1016732), both on the Main track; their scores were not yet available when we submitted this report. Our pre-registered predictions for the final evaluation (Appendix B) will be scored there: P1, that the top-ranked entry's final accuracy lies within two binomial standard errors of the test set's majority constant, and P2, about the entries that led on 7 September. If P1 or P2 fail, the label has learnable structure that we did not find.

\subsection*{J.7\quad Recommendations}

\paragraph{R1. Fix the numerical state and test the fix under load.}
Either configuration of Section 4 made the pipeline byte-reproducible in sequential runs at 5 to 7\% cost; it should be verified under the workers' real load and across cards before it is relied on, and the chosen setting documented, since it selects the fixed point. Isolating submissions on a card may matter as much: two processes on one card were the worst case measured.

\paragraph{R2. Repeat only where it decides something.}
With the re-execution standard deviation measured in the warmup (1.0 item on 17, from four runs), resolving a one-item gap needs about eight runs, though a standard deviation estimated from four runs is itself uncertain, and so is that number; repeating only the entries near a prize boundary at the close costs a handful of runs rather than three times the field.

\paragraph{R3. Publish the labeling prompt, at least after the close.}
For Qwen3.5-4B the answer-format instruction moves the share of problems passing the base condition from 0\% to 88\% (Section 5.3).

\paragraph{R4. Measure the label's reliability directly.}
Two independent draws of variants, labeled with the official process, would give the test-retest agreement that the plug-in simulation cannot (Section 5.2); until then, the reliability of the official label remains unmeasured.

\paragraph{R5. Show the floors and the noise on the leaderboard.}
A row per constant predictor, a stated re-execution standard deviation and binomial intervals; on 21 and 14 items the standard error is 10.9 and 13.4 points.

\paragraph{R6. Require a length control in white-box ablations.}
The White-Box Awards announced on 16 October ask for an ablation showing that the white-box signal matters [6]. On DeepSeek, the organizers' artifact, under its own inference aggregation, ordered the labels mostly through problem length, and our own probe was largely redundant with length as a predictor, adding about two points of AUROC to the character count, with what it reads not established (Section 3); an ablation that removes the white-box component without controlling for length cannot separate the white-box signal from length.

\paragraph{R7. Remove \texttt{permutation\_type} from the aggregated public datasets}
(Section 2).
